\begin{annexitem}{annex:p12.har.prime-logical}{thm:p12-har-prime-logical}{Every selected source logical-library proof}
\begin{proof}
Use \(\mathsf{and}\) projections, the two \(\mathsf{or}\) constructors,
and their generated eliminators with the displayed propositional target.
To discharge the aliases rather than assume a separate logical library,
take and_in to be conj; and_ell and and_elr are the generated and
eliminations into the propositions A and B, with respective branches
\(\lambda u:A.\lambda v:B.u\) and
\(\lambda u:A.\lambda v:B.v\). The no-branch FF eliminator into
a proposition a is false_el. The or constructors are or_inl and
or_inr, and its generated Prop eliminator is or_el. The two predicate
aliases are the actual expressions
\(\mathsf{all\_el}:=\lambda x:A.\lambda h:\Pi y:A.P y.h x\)
and
\(\mathsf{ex\_in}:=\lambda x:A.\lambda u:P x.
\mathsf{ex\_intro}\ A\ P\ x\ u\).
For \(h:\neg(a\vee b)\), the source non_or proof is the pair
\(\langle\lambda u:a.h(\mathsf{or\_inl}\ a\ b\ u),
\lambda v:b.h(\mathsf{or\_inr}\ a\ b\ v)\rangle\).
These operations all conclude propositions and are closed under the
displayed domain/predicate parameters.
Write \(\mathsf{exI}(x,u)\) for the source constructor
\(\mathsf{ex\_intro}\ A\ P\ x\ u\), and
\(\mathsf{exE}_C(e,f)\) for its eliminator into \(C:\mathsf{Prop}\),
where \(f:\Pi x.P x\to C\). Parameters of each instance are the
predicate in that instance, not suppressed changing domains.
The actual four intuitionistic recipes are
\[
\begin{aligned}
 \mathsf{non\_ex}(h)&:=\lambda x.\lambda u:P x.h(\mathsf{exI}(x,u)),\\
 \mathsf{all\_not}(h)&:=\lambda e.\mathsf{exE}_{\mathsf{FF}}
                   (e,\lambda x.\lambda u:P x.h x u),\\
 \mathsf{all\_and}(h)&:=
       \langle\lambda x.\mathsf{and\_ell}(h x),
              \lambda x.\mathsf{and\_elr}(h x)\rangle,\\
 \mathsf{ex\_or}(e)&:=\mathsf{exE}_{(\exists x.P x)\vee(\exists x.Q x)}
       (e,\lambda x.\lambda u:P x\vee Q x.
          \mathsf{orE}(u,
            \lambda v:P x.\mathsf{inl}(\mathsf{exI}(x,v)),
            \lambda w:Q x.\mathsf{inr}(\mathsf{exI}(x,w)))).
\end{aligned}
\]
The fourth instance's input predicate is \(\lambda x.P x\vee Q x\).
The two output constructor instances use \(P\) and \(Q\), respectively.
No witness exists unless supplied by the input constructor, so these
recipes also type check when \(A\) has no closed inhabitant.

For \(\mathsf{dn\_c}\), take
\(\langle\mathsf{DN}\ a,\lambda u:a.\lambda v:\neg a.v u\rangle\).
For \(\mathsf{tnd}\ a\), the explicitly well-typed object is
\[
 \mathsf{DN}(a\vee\neg a)
  (\lambda h:\neg(a\vee\neg a).
      h(\mathsf{inr}(\lambda u:a.h(\mathsf{inl}(u))))).
\]
Here the inner lambda has type \(\neg a\), both calls to \(h\)
have type \(\mathsf{FF}\), and the outer lambda has type
\(\neg\neg(a\vee\neg a)\).
Given \(h:\neg\Pi x.P x\), define
\[
 \mathsf{non\_all}(h):=
 \mathsf{DN}(\exists x.\neg P x)
  (\lambda k:\neg\exists x.\neg P x.
     h(\lambda x.\mathsf{DN}(P x)
         (\lambda u:\neg P x.k(\mathsf{exI}(x,u))))).
\]
Each inner DN proves \(P x\); thus the argument of \(h\) really is
\(\Pi x.P x\). The output predicate of exI is \(\lambda x.\neg P x\).
For \(\mathsf{ex\_c}\), its forward implication is
\(\lambda e.\lambda h:\Pi x.\neg P x.
\mathsf{exE}_{\mathsf{FF}}(e,\lambda x.\lambda u.h x u)\).
Its reverse implication is
\[
 \lambda h:\neg\Pi x.\neg P x.
 \mathsf{DN}(\exists x.P x)
   (\lambda k:\neg\exists x.P x.
      h(\lambda x.\lambda u:P x.k(\mathsf{exI}(x,u)))).
\]
Pair those implications. Every existential/disjunction elimination above
lands in \(\mathsf{Prop}\); DN is not a data-valued choice operation.
\end{proof}
\end{annexitem}

\begin{annexitem}{annex:p12.har.prime-equality}{thm:p12-har-prime-equality}{Equality and separation, with the transport direction checked}
\begin{proof}
Let \(E_{A,x,P}(u,y,e):P y\) denote the \emph{generated} equality
eliminator, with \(u:P x\) and \(e:x=y\); its reflexivity branch is
\(u\). This is the only equality-elimination premise.
For \(e:x=y\), take \(P(z):=(z=x)\) and seed
\(\mathsf{refl\_equal}\ A\ x\) to obtain \(y=x\), giving
\(\mathsf{sym\_eq}\). The source \(\mathsf{leib}\ e\ u\) is precisely
\(E_{A,x,P}(u,y,e)\), in the source's order of arguments.
For \(e:y=x\), first form its just-proved symmetry \(x=y\), then
apply \(E\); this gives \(\mathsf{eq\_ind\_r}\).
For \(f:A\to B\), take \(P(z):=(f x=f z)\) with reflexivity seed;
this yields \(\mathsf{f\_equal}\). Transitivity of equality is the
same transport with \(P(z):=(x=z)\). All these motives are propositions.

Congruence for the data-defined predecessor applied to \(S n=S m\)
gives \(\mathsf{pred}(S n)=\mathsf{pred}(S m)\), definitionally \(n=m\).
For \(e:O=S n\), first obtain \(S n=O\) by symmetry. Transport
\(\mathsf{trivial}:\mathsf{Is\_suc}(S n)\), whose type computes to
\(\mathsf{TT}\), along that
equation, with the \emph{predicate} \(\mathsf{Is\_suc}\); the result
has type \(\mathsf{Is\_suc}(O)=\mathsf{FF}\). Thus
\(\mathsf{no\_conf}\ n:=\lambda e.E(\mathsf{trivial},O,
\mathsf{sym\_eq}\ e)\). Reversing the equality without reversing
the seed would not give this contradiction. The converse separation
\(\neg(S n=O)\) follows by symmetry. Ex falso is used only to
prove another proposition, never to select natural-number data.
\end{proof}
\end{annexitem}

\begin{annexitem}{annex:p12.har.prime-order}{thm:p12-har-prime-order}{All order inversions and the actual course-of-values construction}
\begin{proof}
Every induction mentioned here has an explicit \(\mathsf{Prop}\)
motive. First \(O\le n\) follows by nat induction: reflexivity at
\(O\), and \(\mathsf{leseq\_suc}\) in the successor step.
For \(x\le y\le z\), induct on the second \(\mathsf{leseq}\) proof
with fixed \(y\) and motive \(\lambda z.\,(x\le y)\to(x\le z)\).
Its reflexivity branch returns the input proof of \(x\le y\), and
its successor branch applies \(\mathsf{leseq\_suc}\) to the induction
result. This is \(\mathsf{leseq\_trans}\).
Induction on \(n\le m\) with motive \(S n\le S m\) gives successor
monotonicity: its reflexivity branch is reflexivity at \(S n\), and
its successor branch again uses \(\mathsf{leseq\_suc}\).
Also \(n\le S n\) is the constructor applied to reflexivity.

For \(S n\le m\), induction on this proof with motive
\(\mathsf{Is\_suc}(m)\) has seed \(\mathsf{trivial}\) and constant
successor branch \(\mathsf{trivial}\). At \(m=O\) this is
\(\mathsf{FF}\), proving \(\mathsf{non\_lt0}\).
The data case split on \(m\) gives \(\mathsf{pred}(m)\le m\):
reflexivity if \(m=O\), and \(k\le S k\) if \(m=S k\).
Now induct on \(S n\le m\) with motive \(n\le\mathsf{pred}(m)\).
At \(m=S n\) use reflexivity. At the step from \(m=k\) to \(S k\),
the induction hypothesis is \(n\le\mathsf{pred}(k)\);
compose it with \(\mathsf{pred}(k)\le k\), which is the required
\(n\le\mathsf{pred}(S k)\). Specializing \(m=S r\) proves
\(\mathsf{suc\_leseq}\). Thus
\(n<S m\leftrightarrow n\le m\): one direction is successor
cancellation and the other is successor monotonicity.
Compose \(n\le S n\) with \(S n\le m\) for
\(\mathsf{lthan\_leseq}\).

Strict irreflexivity \(\neg(n<n)\) follows by nat induction. At
zero it is \(\mathsf{non\_lt0}\). At \(S n\), a putative
\(S(S n)\le S n\) cancels to \(S n\le n\), contradicting the
induction hypothesis. A strict/weak chain therefore cannot cycle:
\(a<b\) and \(b\le a\) compose to \(S a\le a\).
Construct trichotomy with motive
\(\Pi m.\,(n=m)\vee(n<m)\vee(m<n)\), by induction on \(n\)
and data case analysis on \(m\). At \(n=m=O\) give reflexivity.
At \(n=O,m=S k\), lift \(O\le k\) to \(S O\le S k\).
At \(n=S k,m=O\), give the opposite strict inequality similarly.
At \(n=S k,m=S l\), use the induction hypothesis for \(k,l\):
lift equality by congruence, and each strict inequality by successor
monotonicity. These exhaust the constructor cases, without DN.
Antisymmetry follows by that trichotomy and the just-proved impossibility
of a strict/weak cycle. Zero minimality follows by cases on \(n\):
\(S k\le O\) is \(k<O\), excluded by non_lt0.

For \(\mathsf{lt01}\), split \(x\): zero is the left alternative;
\(S O\) is the middle alternative; \(S(S k)\) satisfies
\(\mathsf{one}<S(S k)\), by lifting \(O\le k\) twice.
Eliminating the first two alternatives against the supplied negative
equations gives \(\mathsf{n0n1lt}\).
If \(k\le S n\), apply trichotomy to \(k,S n\). Equality is one
alternative; \(k<S n\) cancels to \(k\le n\); and
\(S n<k\) contradicts \(k\le S n\) by the strict/weak cycle.
This proves the successor-bound split used below.

For the course-of-values theorem, let
\(H:\Pi n.\mathsf{before}\ n\ P\to P n\) and use ordinary nat
induction on the \emph{proposition}
\(Q n:=\Pi k.(k\le n)\to P k\). At zero, prove \(P O\)
by \(H O\) applied to \(\lambda k.\lambda u:k<O.
\mathsf{false\_el}(\mathsf{non\_lt0}\ k\ u)\).
For \(k\le O\), zero minimality gives \(k=O\); reverse this
equation and transport \(P O\) to \(P k\), supplying \(Q O\).
At \(S n\), the induction hypothesis \(Q n\) gives every
\(P k\) for \(k<S n\), by the equivalence \(k<S n\iff k\le n\).
Apply \(H(S n)\) to this function to obtain \(P(S n)\).
For arbitrary \(k\le S n\), the successor-bound split gives
either the old \(Q n\) case or an equality \(k=S n\), along whose
symmetry the new \(P(S n)\) is transported. This constructs
\(Q(S n)\). Finally \(Q n\ n\ \mathsf{leseq\_n}\) proves \(P n\).
All displayed functions are proofs of propositions; no function uses
an ordering proof to return a data-valued bound or predecessor.
\end{proof}
\end{annexitem}

\begin{annexitem}{annex:p12.har.prime-semiring}{thm:p12-har-prime-semiring}{Every source arithmetic identity and the positive-factor order argument}
\begin{proof}
All equalities here are the source propositional equality on nat. Each
chain composes the equality transport/transitivity just constructed;
replacing a subexpression uses f_equal, not equality reflection.
Computational endpoints use the displayed structural-recursion equations.
The induction motives are the displayed equations universally quantified
over the remaining natural-number parameters.

For \(\mathsf{plus\_altsuc}\), induct on \(n\).
At \(O\), both sides compute to \(S m\). At \(S n\), they are
\(S(n+S m)\) and \(S(S(n+m))\), so apply congruence for S to the
induction hypothesis. For \(\mathsf{plus\_altzero}\), zero computes
to zero; its successor step is congruence for S on \(n+O=n\).
For \(\mathsf{plus\_ass}:n+(m+k)=(n+m)+k\), zero computes to
\(m+k\); the successor step computes both sides to S applied to
the previous equality. For \(\mathsf{plus\_com}\), at zero use
the symmetry of \(m+O=m\). At \(S n\) the actual chain is
\[
 (S n)+m=S(n+m)=S(m+n)=m+(S n),
\]
using the induction hypothesis and the symmetry of plus_altsuc.
These proofs establish, as a usable equality recipe,
\((a+b)+c=(a+c)+b\): rewrite the left side to \(a+(b+c)\),
commute \(b+c\), then reassociate. They also establish arbitrary
reassociation and permutation of a finite sum by repeated applications
of the same two proved laws; no additional arithmetic axiom is meant.

Addition cancellation \(a+b=a+c\to b=c\) is nat induction on a.
At zero the premise computes to the conclusion. At S a, successor
injectivity first yields \(a+b=a+c\), then the induction hypothesis
applies. Right cancellation follows by commuting each side and applying
this result. In particular \(n+\mathsf{one}=S n\), by plus_altsuc
at O and plus_altzero.

For \(\mathsf{timesaltzero}\), induction on n has zero reflexivity;
the step is \((S n)O=nO+O=O+O=O\).
For \(\mathsf{timesaltsuc}:n(S m)=nm+n\), induction on n again
has the zero computation. Set \(a=nm\) in the successor step.
Its left side is
\[
 (S n)(S m)=n(S m)+S m=(a+n)+S m=S((a+n)+m),
\]
and its right side is
\[
 ((S n)m)+S n=(a+m)+S n=S((a+m)+n).
\]
The finite-sum recipe just proved gives
\((a+n)+m=(a+m)+n\), completing the step by S congruence.
For \(\mathsf{distr}:(n+m)k=nk+mk\), induction on n has
\(mk=O+mk\) at zero. The successor step is the concrete chain
\[
 ((S n)+m)k=(S(n+m))k=(n+m)k+k
       =(nk+mk)+k=(nk+k)+mk=((S n)k)+mk.
\]
For \(\mathsf{times\_ass}:n(mk)=(nm)k\), induction on n has
zero computation. Its step is
\[
 (S n)(mk)=n(mk)+mk=(nm)k+mk=(nm+m)k=((S n)m)k,
\]
where the penultimate equation is the symmetry of distr.
Multiplication commutativity is a separate derived fact, by induction
on n: zero uses timesaltzero at m, and
\[
 (S n)m=nm+m=mn+m=m(S n)
\]
uses the induction hypothesis and timesaltsuc in the last step.
Right distributivity now follows without a new premise:
\(n(m+k)=(m+k)n=mn+kn=nm+nk\).
Also \(\mathsf{one}\,n=On+n=O+n=n\); commute for
\(n\,\mathsf{one}=n\).

For the additive characterization of order, induct on a proof of
\(n\le m\), eliminating into \(\exists k.n+k=m:\mathsf{Prop}\).
The reflexivity branch is witness O, with plus_altzero.
At a successor branch, eliminate the induction result into the same
propositional existential: from local witness k and equation
\(n+k=m\), introduce witness \(S k\), since
\(n+S k=S(n+k)=S m\).
For the converse, induction on k constructs \(n\le n+k\): at
zero transport reflexivity along plus_altzero; at the successor use
the leseq successor constructor and \(n+S k=S(n+k)\).
Then eliminate \(\exists k.n+k=m\) into \(n\le m\) and transport
this bound along the local equation. This also proves
\(n\le n+k\), and by commutation \(k\le n+k\).

To make the positive-factor reflection argument explicit, first prove
\(a<a+S u\). The already constructed \(a\le a+u\) lifts to
\(S a\le S(a+u)=a+S u\), which is the required strict inequality.
Any \(d\) with \(O<d\) has a successor constructor shape: a data
case split on d excludes O by non_lt0, leaving \(d=S u\).
Consequently \((S k)d=kd+d=kd+S u=S(kd+u)\) is positive,
and \(a<a+(S k)d\), by the preceding strict-sum fact and this
equation. If \(r<s\), the additive characterization for
\(S r\le s\) supplies, locally in Prop, k with
\(S r+k=s\). The addition equations give \(s=r+S k\), and distr
gives \(sd=rd+(S k)d\). The strict-sum argument gives \(rd<sd\).
Thus positive multiplication preserves strict order.
If \(rd\le sd\), trichotomy of r,s gives equality, \(r<s\), or
\(s<r\). The first two imply \(r\le s\), respectively by
transported reflexivity and lthan_leseq. The last implies \(sd<rd\),
whose composition with \(rd\le sd\) contradicts strict irreflexivity
at sd. This proves the asserted weak-order reflection, without DN.
For later use, a positive multiple \((S k)d\) is at least d,
because \((S k)d=kd+d=d+kd\) and \(d\le d+kd\).
\end{proof}
\end{annexitem}

\begin{annexitem}{annex:p12.har.prime-divisibility}{thm:p12-har-prime-divisibility}{Local coefficient arithmetic, including the zero-divisor case}
\begin{proof}
Every existential elimination in this proof has a propositional target.
For reflexivity \(\mathsf{div}\ n\ n\), introduce coefficient one,
since \(\mathsf{one}\,n=n\). For \(\mathsf{div}\ \mathsf{one}\ n\),
introduce coefficient n, since \(n\,\mathsf{one}=n\).
For transitivity, eliminate \(\mathsf{div}\ d\ m\) and
\(\mathsf{div}\ m\ n\) locally: let \(ad=m\), \(bm=n\).
Associativity and substitution give \((ba)d=b(ad)=bm=n\).
Introduce ba as the coefficient in \(\mathsf{div}\ d\ n\).

Suppose \(\mathsf{one}<d\) and \(rd=\mathsf{one}\).
Split the \emph{data} r. For \(r=O\), the equation computes to
\(O=S O\), contradicted by no_conf. For \(r=S k\), the proved
positive-multiple bound gives \(d\le rd=\mathsf{one}\).
Composing \(S\mathsf{one}\le d\) with that bound gives
\(\mathsf{one}<\mathsf{one}\), again a contradiction. Eliminating
the divisibility existential into FF proves that d cannot divide one.

The subtraction equations give \(a\mathbin{\dot{-}}O=a\), by
cases on a. Nat induction on r with all k quantified proves
\[
 (r+k)\mathbin{\dot{-}}r=k.
\]
At r=O, the result is \(k\mathbin{\dot{-}}O=k\). At S r,
\(((S r)+k)\mathbin{\dot{-}}(S r)\) computes to
\((r+k)\mathbin{\dot{-}}r\), so the induction hypothesis applies.
Thus if \(r+k=s\), propositional congruence and symmetry give
\(s\mathbin{\dot{-}}r=k\). This is an equation about already
defined natural-number data, not an extraction from an ordering proof.

Now let \(\mathsf{div}\ d\ a\) and \(\mathsf{div}\ d\ (a+b)\).
Eliminate these two existentials into the target
\(\mathsf{div}\ d\ b:\mathsf{Prop}\), retaining local r,s with
\(rd=a\), \(sd=a+b\). Split the data d.
If \(d=O\), timesaltzero yields \(a=O\) and \(a+b=O\).
Substitute \(a=O\); then \(O+b=b\) gives \(b=O\).
Introduce coefficient O in \(\mathsf{div}\ O\ b\), transporting
the reflexive computation \(OO=O\) to b. This settles the case
which cancellation by a zero multiplier cannot settle.

If \(d=S u\), then \(O<d\), by successor monotonicity of
\(O\le u\). The additive order fact \(a\le a+b\) transports
along the two local equations to \(rd\le sd\). Positive-factor
reflection therefore gives \(r\le s\). Eliminate the additive
characterization into the still propositional target: take local k
with \(r+k=s\). Distr gives
\[
 a+b=sd=(r+k)d=rd+kd=a+kd.
\]
Addition cancellation yields \(b=kd\). The subtraction lemma gives
\(s\mathbin{\dot{-}}r=k\), so substitution gives
\((s\mathbin{\dot{-}}r)d=b\). Introduce the explicit local
coefficient \(s\mathbin{\dot{-}}r\) in the output existential.
The two coefficient inputs came from ex eliminations into Prop; the
output is another ex proof in Prop. This derivation cannot be read as
a map \(\mathsf{div}\ d\ n\to\mathsf{nat}\) or as classical choice.
\end{proof}
\end{annexitem}

\begin{annexitem}{annex:p12.har.prime-divisor}{thm:p12-har-prime-divisor}{The complete course-of-values prime-divisor proof}
\begin{proof}
Set \(P n:=(\mathsf{one}<n)\to\mathsf{HPD}\ n\), which is a
well-formed \(\mathsf{Prop}\)-valued predicate on the source nat.
The proved cv_ind reduces the target to construction of
\(\Pi n.(\Pi d.(d<n)\to P d)\to P n\).
Take arbitrary n, \(I:\Pi d.(d<n)\to P d\), and
\(h:\mathsf{one}<n\). Instantiate the explicitly derived tnd at
\(B_n:=\exists d.\mathsf{propdiv}\ d\ n\). Its or elimination
target is \(\mathsf{HPD}\ n:\mathsf{Prop}\).

In its negative branch \(u:\neg B_n\), construct
\(\langle h,u\rangle:\mathsf{prime}\ n\), and use divisibility
reflexivity \(\mathsf{div}\ n\ n\). The ex constructor with
witness n and this pair gives \(\mathsf{HPD}\ n\).
In its positive branch \(e:B_n\), use ex elimination into that
same target, with arbitrary local d and
\(v:(\mathsf{one}<d)\wedge(d<n)\wedge\mathsf{div}\ d\ n\).
Project its three conjuncts. The induction input is exactly
\(I\ d\ (d<n)\ (\mathsf{one}<d):\mathsf{HPD}\ d\).
Eliminate this existential, again into \(\mathsf{HPD}\ n\), with
local p and pair \(w:\mathsf{prime}\ p\wedge\mathsf{div}\ p\ d\).
Divisibility transitivity applied to its second projection and the
third projection of v gives \(\mathsf{div}\ p\ n\).
Pair this with the first projection of w and introduce p into
\(\mathsf{HPD}\ n\). This completes both branches and hence the
step required by cv_ind. Abstract n for the stated declaration.

The sole classical instance in the step is tnd at \(B_n\), whose
DN recipe is given in the logical-library annex; neither the induction
principle nor any arithmetic lemma here uses DN. The development proves
inhabitation of a Prop existential, not a Set-valued algorithm finding
its prime. A constructive search would require its own terminating
data procedure and correctness proof, not the elimination used here.
\end{proof}
\end{annexitem}

\begin{annexitem}{annex:p12.har.prime-euclid}{thm:p12-har-prime-euclid}{Factorial divisibility and the fully linked Euclid argument}
\begin{proof}
First construct \(\mathsf{one}\le\mathsf{fac}\ n\) by nat
induction. At zero this is reflexivity, since fac O=one.
At S n, the factorial equation and the multiplication successor
equation give
\[
 \mathsf{fac}(S n)=(S n)(\mathsf{fac}\ n)
       =n(\mathsf{fac}\ n)+\mathsf{fac}\ n
       =\mathsf{fac}\ n+n(\mathsf{fac}\ n).
\]
The additive bound gives \(\mathsf{fac}\ n\le\mathsf{fac}(S n)\).
Compose with the induction hypothesis. This is a proof of a bound,
not a property inferred just from the recursive definition.

Next use nat induction on n with the quantified propositional motive
\[
 F n:=\Pi m.(\mathsf{one}\le m)\to(m\le n)\to
                 \mathsf{div}\ m\ (\mathsf{fac}\ n).
\]
At zero, \(\mathsf{one}\le m\le O\) composes to
\(\mathsf{one}\le O\), which is excluded by non_lt0 at O;
ex falso proves the required divisibility proposition.
At S n, fix m and its two bounds. The already proved successor-bound
split on \(m\le S n\) gives \(m\le n\) or \(m=S n\).
If \(m=S n\), introduce coefficient \(\mathsf{fac}\ n\), since
commutativity gives
\((\mathsf{fac}\ n)(S n)=(S n)(\mathsf{fac}\ n)
=\mathsf{fac}(S n)\), and transport along the equation for m.
If \(m\le n\), apply \(F n\ m\) to its bounds, yielding
\(\mathsf{div}\ m\ (\mathsf{fac}\ n)\). Eliminate this
existential into the desired divisibility proposition: from local
coefficient a with \(am=\mathsf{fac}\ n\), associativity gives
\[
 ((S n)a)m=(S n)(am)=(S n)(\mathsf{fac}\ n)
                          =\mathsf{fac}(S n).
\]
Introduce coefficient \((S n)a\). This closes both branches of the
inductive step and proves factorial divisibility at every n.

Fix arbitrary N, and form the data
\(M:=\mathsf{fac}\ N+\mathsf{one}=S(\mathsf{fac}\ N)\).
The bound just proved is \(S O\le\mathsf{fac}\ N\), or
\(O<\mathsf{fac}\ N\). Successor monotonicity therefore gives
\(S(S O)\le S(\mathsf{fac}\ N)=M\), exactly
\(\mathsf{one}<M\). Apply numbers_gt1_have_primediv to obtain
\(\mathsf{HPD}\ M\). Eliminate it into the propositional target
\(\exists p.\mathsf{prime}\ p\wedge N<p\), with local p and
\(u:\mathsf{prime}\ p\wedge\mathsf{div}\ p\ M\).
Project \(\mathsf{one}<p\) from its primality component; in
particular \(\mathsf{one}\le p\) by lthan_leseq.

If \(p\le N\), factorial divisibility yields
\(\mathsf{div}\ p\ (\mathsf{fac}\ N)\). Apply the proved
divisibility subtraction theorem with
\(a=\mathsf{fac}\ N\), \(b=\mathsf{one}\), and the given
\(\mathsf{div}\ p\ M\); this yields
\(\mathsf{div}\ p\ \mathsf{one}\), contradicted by
\(\mathsf{one}<p\) and the prime-divisibility library.
Thus \(\neg(p\le N)\). The constructive trichotomy of N,p now
has only the alternative \(N<p\): equality would transport
reflexivity to \(p\le N\), and \(p<N\) would give the same
bound by lthan_leseq. Pair \(N<p\) with the given primality of p,
and introduce that local p into the target existential. Abstract N.

All witnesses used in factorial divisibility and the final prime
statement were introduced by ex constructors or locally bound by ex
eliminations into Prop. The factorial/subtraction functions themselves
are data recursions independent of those proofs. DN enters this theorem
only through the explicitly linked prime-divisor theorem; there is no
extra appeal to classical choice, modern library automation or a
different definition of primality.

For the finite-family formulation, construct
\(B k\ a:=\mathsf{nat\_rec}\ O\
(\lambda j:\mathsf{nat}.\lambda r:\mathsf{nat}.r+a j)\ k\),
an instance with data target nat and the indicated parameter a.
Its equations are the displayed fold equations. Nat induction on k
with propositional motive
\(\Pi i.(i<k)\to(a i\le B k\ a)\) proves the bound.
At zero the premise is excluded by non_lt0.
At S k, cancel \(i<S k\) to \(i\le k\) and use trichotomy of
i,k. In the equality branch, \(a i=a k\) by f_equal and
\(a k\le B k\ a+a k=B(S k)\ a\) by the additive bound.
In the \(i<k\) branch, the induction hypothesis gives
\(a i\le B k\ a\); compose with
\(B k\ a\le B k\ a+a k\). The \(k<i\) branch contradicts
\(i\le k\). This proves every required bound.
Apply the already proved Euclid result at the data bound \(B k\ a\)
and eliminate its existential into the finite-family existential
proposition. For local prime p with \(B k\ a<p\), take arbitrary
\(i<k\). A putative \(p=a i\) transports the proved
\(a i\le B k\ a\) to \(p\le B k\ a\), a strict/weak cycle.
For an explicit proof recipe, name the proved bound function H,
put \(T:=B k\ a\), and let \(u_{>}:T<p\) be the local Euclid
proof. Let \(\mathsf{irr}_T:(T<T)\to\mathsf{FF}\) be the constructed
strict-irreflexivity proof. The required negative-equation function is
\[
\begin{aligned}
 &\lambda i:\mathsf{nat}.\lambda h:i<k.\lambda e:p=a i.\\
 &\quad\mathsf{irr}_T\bigl(
  \mathsf{leseq\_trans}\ (S T)\ p\ T\ u_{>}
   (\mathsf{leib}\ \mathsf{nat}\ (\lambda z:\mathsf{nat}.z\le T)
       \ (a i)\ p\ (\mathsf{sym\_eq}\ e)\ (H i h))\bigr).
\end{aligned}
\]
The symmetry argument has type \(a i=p\), so leib really returns
\(p\le T\); transitivity then returns \(S T\le T\), the input
type of \(\mathsf{irr}_T\). Pair the function with primality and
introduce p. All resulting quantifiers
and existential eliminations are still in Prop. No List type, sequence
code, maximum operation, or additional recursion theorem was imported.
\end{proof}
\end{annexitem}
